\chapter*{How to Use This Book}
\addcontentsline{toc}{chapter}{How to Use This Book}
\markboth{How to Use This Book}{How to Use This Book}
This book is a complete, code-level explanation of the \emph{doc-agent} repository: an \textbf{agentic retrieval-augmented generation (RAG)} system that answers questions about \emph{scanned documents}. Every function that you will be asked about in the exam is shown here, \emph{verbatim from the source}, followed by an explanation of what it does, why it exists, and what could go wrong.

\section*{What the exam will ask}
The final exam is solved \emph{using the codebase}. Expect three kinds of questions:
\begin{enumerate}
 \item \textbf{Understanding}: ``What does the hook seam \c{before\_answer} do, and which modules attach to it?''
 \item \textbf{Tracing / computing by hand}: ``Given this config and this retrieval score, how many retrieval steps does the agent take, and what is the answer?''
 \item \textbf{Modifying / extending}: ``Add a new PII pattern'', ``Add a tool'', ``Change the policy so that \dots'', ``This test fails; find the bug.''
\end{enumerate}
Chapter~\ref{ch:problems} contains worked problems of all three kinds. Try them \emph{before} reading the solutions.

\section*{Reading order}
\begin{center}
\begin{tabular}{@{}rll@{}}
\toprule
Chapter & Topic & Read it when \\
\midrule
1 & Big picture, vocabulary & first, always \\
2 & Repo layout and the ``rules'' (what is FIXED) & first \\
3 & Data contracts + configuration & before any stage \\
4 & \textbf{Hooks and wiring} (the cross-cutting design) & before the agent chapter \\
5--7 & Phase A: building the knowledge base & in order \\
8 & Phase B: the agent (most examined) & carefully, twice \\
9 & Cross-cutting features end to end & after 4 and 8 \\
10--12 & Evaluation, RL, training, MLOps, serving, tests & reference \\
13 & Design decisions, trade-offs, limits & for essay questions \\
14 & Exam-style problems with solutions & last, and repeatedly \\
\bottomrule
\end{tabular}
\end{center}

\section*{Conventions}
\begin{itemize}
 \item Code listings are extracted automatically from the repository, so they match the code exactly. Line numbers inside a listing are relative to that excerpt.
 \item \texttt{FIXED} means the file or function was supplied and \emph{must not} change (the structure gate in CI checks this). \texttt{IMPL} means it was a stub that has been implemented.
 \item Coloured boxes: \textbf{Why it exists} (design reason), \textbf{Gotcha} (common mistake / subtle behaviour), \textbf{Exam tip}, \textbf{Key idea}, \textbf{Try it} (experiment).
\end{itemize}

\begin{keyidea}
If you only remember five things:
\begin{enumerate}
 \item The system has \textbf{two phases}: build the index once (Stages 1--4), answer questions many times (Stages 5--6).
 \item The \textbf{agent} is a loop whose path depends on a number it observes: the best retrieval score (\emph{evidence-gated re-search}).
 \item Safety, privacy, tracing and grounding are \textbf{not} inside the stages; they plug into fixed \textbf{hook seams}, registered in one place (\c{wiring.py}).
 \item The system must \textbf{abstain} rather than hallucinate: no strong evidence means ``insufficient evidence''.
 \item All models and numbers live in \textbf{\c{configs/config.yaml}}; code reads from it, never hard-codes.
\end{enumerate}
\end{keyidea}
